\documentclass[12pt,a4paper]{article}
\usepackage{style/main}
\usepackage{mathtools}
\usepackage{tikz}
\usetikzlibrary{decorations.pathreplacing}

\startdocument{Chimie}

\titre{Chapitre 2 - Méthodes physiques d'analyse}{Spectroscopie et conductimétrie}

\section*{I - Rappels: Spectre et couleurs}
\subsection*{A - Rappels de secondes}
La couleur perçue correspond à la partie de la lumière \textbf{réfléchie ou transmise}, c'est-à-dire celle qui n'est pas absorbée.  
\\Le spectre visible s'étend de \textbf{400 nm (violet)} à \textbf{800 nm (rouge)}.
\subsection*{B - Nomenclature}
\vspace{3\baselineskip}
\subsection*{C - Lien entre couleur et structure moléculaire}

Plus une molécule comporte de \textbf{liaisons doubles conjuguées}, plus elle absorbe à \textbf{grande longueur d'onde}.  
\\La présence de certains \textbf{groupes} (OH, NH$_2$, NO$_2$, ou certains halogènes) influe le domaine d'absorption des molécules.
\vspace{1\baselineskip}
\section*{II - Spectroscopie UV-Visible}
\subsection*{A - Abosrbance et Loi de Beer-Lambert}

\formule{
L'absorbance $A$ est reliée à la concentration par la loi de Beer-Lambert :

\[
A =
\begin{tikzpicture}[baseline=(current bounding box.center)]
\node (vals) {
$\begin{array}{l}
\varepsilon_\lambda \\
l \\
c
\end{array}$
};
\draw[decorate,decoration={brace,mirror,amplitude=7pt}] 
(vals.north west) -- (vals.south west) node[midway,xshift=-1.2cm] {Paramètres};
\node[right=1.5cm of vals.east, anchor=west] {
$\begin{array}{l}
\text{coeff d'extinction molaire (L·mol$^{-1}$·cm$^{-1}$)} \\
\text{largeur de la cuve (cm)} \\
\text{concentration molaire (mol·L$^{-1}$)}
\end{array}$
};
\end{tikzpicture}
\]
}

\subsection*{B - Spectrophotométrie}
La spectrophotométrie est la mesure du spectre d'absorbance à l'aide d'un spectrophotomètre.  
La grandeur mesurée est \textbf{l'absorbance A}.

\subsection*{C - Facteurs influençant A}
L'absorbance dépend :  
– du coefficient $\varepsilon(\lambda)$,  
– de la concentration $c$,  
– de la largeur de cuve $l$,  
– de l'espèce chimique étudiée.

\subsection*{D - Dosage spectrophotométrique}
\exemple{
Procédure :  
1. Préparer une gamme étalon.  
2. Mesurer $A = f(\lambda)$.  
3. Régler sur $\lambda_{max}$ et faire le blanc.  
4. Mesurer $A$ des solutions étalons et de l'inconnue.  
5. Tracer $A = f(c)$ et interpoler la concentration inconnue.  
}

---

\section*{III - Conductivité des solutions électrolytiques}

\subsection*{A - Loi d'Ohm}
\formule{
\[
U = R \cdot I
\]
avec $U$ la tension (V), $R$ la résistance ($\Omega$), $I$ l'intensité (A).
}

\subsection*{B - Conductance}
\formule{
\[
G = \frac{1}{R} = \frac{I}{U}
\]
}
Unité : Siemens (S).

\subsection*{C - Facteurs influençant $G$}
La conductance dépend :  
– de la nature des ions,  
– de la concentration (augmente avec $c$),  
– de la surface des électrodes,  
– de la température (mobilité ionique),  
– de la distance entre électrodes.

\subsection*{D - Conductivité}
\formule{
\[
G = k \cdot \sigma
\]
}
où $\sigma$ est la conductivité (S·m$^{-1}$) et $k$ la constante de cellule.

\subsection*{E - Conductivité molaire ionique}
\formule{
\[
\sigma = \sum_i \lambda_i \cdot [X_i]
\]
}
Chaque ion possède une conductivité molaire ionique propre $\lambda_i$ (loi de Kohlrausch).

\subsection*{F - Dosage conductimétrique}
\exemple{
Procédure :  
1. Préparer une gamme étalon.  
2. Mesurer $\sigma$ des solutions.  
3. Tracer $\sigma = f(c)$.  
4. Déterminer la concentration de l'inconnue par interpolation.  
}

\subsection*{G - Cas particulier : dosage d'un gaz}
\formule{
\[
PV = nRT
\]
}
À $V$, $R$ et $T$ constants, la pression est proportionnelle à la quantité de matière dissoute.

---

\section*{IV - Spectroscopie infrarouge (IR)}

\subsection*{A - Groupes caractéristiques}
\formule{
\[
E = h \cdot \nu = \frac{h \cdot c}{\lambda}
\]
}
Groupes observables : carboxyle, hydroxyle, carbonyle, amine, ester, alcène…

\subsection*{B - Présentation du spectre}
Le spectre IR se trace en transmittance en fonction du nombre d'onde $\tilde{\nu}$ (cm$^{-1}$), entre 4000 et 500 cm$^{-1}$.  
À gauche : grandes énergies (liaisons fortes).  
À droite : faibles énergies (infrarouge lointain).

\subsection*{C - Origine des bandes IR}
Les atomes vibrent par :  
– élongation (liaison qui s'allonge/compresse),  
– déformation angulaire.  
Quand l'énergie du rayonnement correspond à celle de
\end{document}
